\section{Lab 2: Kinematic Modeling using MATLAB Simulink}

\subsection{Objective}
To understand basic kinematic models of mobile robots (bicycle model and unicycle model) and implement them using MATLAB Simulink.

\subsection{Requirements}
\begin{itemize}
	\item \textbf{Hardware:} Computer with minimum 4GB RAM
	\item \textbf{Software:} 
	\begin{itemize}
		\item MATLAB R2020a or later
		\item Simulink
	\end{itemize}
\end{itemize}

\subsection{Theory}

\subsubsection{Unicycle Model}
The unicycle model is a simple way to describe robot motion. It uses linear velocity (v) and angular velocity (w). The equations are:

\begin{equation}
\dot{x} = v \cos(\theta)
\end{equation}
\begin{equation}
\dot{y} = v \sin(\theta)
\end{equation}
\begin{equation}
\dot{\theta} = \omega
\end{equation}

where (x, y) is position and theta is orientation angle.

\subsubsection{Bicycle Model}
The bicycle model is used for car-like robots. It considers wheelbase L and steering angle. The equations are:

\begin{equation}
\dot{x} = v \cos(\theta)
\end{equation}
\begin{equation}
\dot{y} = v \sin(\theta)
\end{equation}
\begin{equation}
\dot{\theta} = \frac{v}{L} \tan(\delta)
\end{equation}

where v is velocity, L is wheelbase distance, and delta is steering angle.

\subsection{Procedure}

\subsubsection{Creating Unicycle Model}

\begin{enumerate}
	\item Open MATLAB and type \texttt{simulink} in command window
	\item Create new model
	\item Add these blocks:
	\begin{itemize}
		\item Constant blocks for v and omega
		\item Integrator blocks for x, y, and theta
		\item Trigonometric blocks (sin, cos)
		\item Product blocks for multiplication
		\item XY Graph for plotting path
	\end{itemize}
	\item Connect blocks according to equations
	\item Set values: v = 1 m/s, omega = 0.5 rad/s
	\item Run simulation for 20 seconds
\end{enumerate}

\subsubsection{Creating Bicycle Model}

\begin{enumerate}
	\item Create new Simulink model
	\item Add blocks similar to unicycle model
	\item Additional blocks needed:
	\begin{itemize}
		\item Constant for wheelbase L
		\item Constant for steering angle
		\item Trigonometric block for tan
		\item Gain block for 1/L
	\end{itemize}
	\item Connect according to bicycle equations
	\item Set values: v = 2 m/s, L = 2.5 m, steering = 15 degrees
	\item Run simulation
\end{enumerate}

\subsubsection{Block Diagram}
% Add your Simulink screenshot here
\begin{figure}[h]
	\centering
	% \includegraphics[width=0.8\textwidth]{images/bicycle_model.png}
	\caption{Simulink block diagram (Add your image here)}
	\label{fig:model}
\end{figure}

\subsubsection{Code}
MATLAB initialization script:
\begin{lstlisting}[language=Matlab]
% Parameters for unicycle model
v_uni = 1;      % velocity m/s
omega = 0.5;    % angular velocity rad/s

% Parameters for bicycle model
v_bike = 2;     % velocity m/s
L = 2.5;        % wheelbase m
delta = 15;     % steering angle degrees
delta_rad = delta * pi/180;

% Simulation time
t_sim = 20;

% Plot results after simulation
figure;
plot(x_uni, y_uni, 'b-', 'LineWidth', 2);
hold on;
plot(x_bike, y_bike, 'r-', 'LineWidth', 2);
grid on;
xlabel('X Position (m)');
ylabel('Y Position (m)');
legend('Unicycle', 'Bicycle');
title('Robot Trajectories');
\end{lstlisting}

\subsection{Output}

The simulations showed different paths:

\textbf{Unicycle Model:} Made a circular path because we gave constant velocities. The radius was around 2 meters.

\textbf{Bicycle Model:} Made a curved path with a bigger turning radius than the unicycle. This is because of the wheelbase constraint.

Both models moved smoothly and the paths looked correct based on the equations we used.

\subsection{Inference}

What we learned from the outputs:
\begin{itemize}
	\item Unicycle model makes circles when you give constant inputs
	\item Bicycle model turns wider because of the wheelbase length
	\item Changing velocity affects how fast the robot moves but not the path shape
	\item Changing angular velocity or steering angle changes the turning radius
	\item Unicycle model is simpler - good for two-wheeled robots
	\item Bicycle model is better for car-like vehicles
	\item Simulink makes it easy to see how equations translate to actual motion
\end{itemize}

\subsection{Conclusions}

We learned how to model robot motion using kinematic equations in Simulink. The unicycle model is simple and works for differential drive robots. The bicycle model is more realistic for car-like robots because it includes wheelbase and steering constraints. Both models help understand how robots move before building actual hardware. The visual block diagram in Simulink made it easier to understand the math behind robot motion.

\clearpage